%!TEX TS-program = xelatex
\documentclass[]{project-reference}
\usepackage{afterpage}
\usepackage{hyperref}
\usepackage{xcolor}

\hypersetup{
    pdftitle={Project Reference - Maximilian Zuleger},
    pdfauthor={Maximilian Zuleger},
    pdfsubject={Project Reference},
    pdfkeywords={cv,maxi,zuleger,software,projects},
    colorlinks=false,
    allbordercolors=white
}

\begin{document}
\header{Maximilian~}{Zuleger}
      {Project References}
\hspace{5cm}
\vspace{0.4cm}

\section{Projects}

% ---------------------------------------------------------------------------
% EXAMPLE ENTRY (copy this) - please replace/extend with real projects.
% Just duplicate this block per project and fill in the details; all
% styling comes automatically from project-reference.cls.
% ---------------------------------------------------------------------------

\projectentry{Example Client GmbH -- Order Management Platform}{Replaced a legacy monolith with a scalable microservices platform for B2B order processing.}

\reflabel{Tech-Stack}
\techstack{Java, Spring Boot, Kafka, PostgreSQL, Docker, Kubernetes, AWS}

\reflabel{Highlights}
\begin{itemize}
  \item Led a 5-person engineering team, introduced Scrum and CI/CD practices
  \item Designed and implemented the microservices architecture for B2B order processing and logistics
  \item Integrated with SAP systems and coordinated requirements together with business stakeholders
  \item Mentored junior developers and conducted code reviews
\end{itemize}

\projectsep

% Add the next project entry here, e.g.:
%
% \projectentry{Company XY -- Short project title}{One sentence describing the project.}
%
% \reflabel{Tech-Stack}
% \techstack{...}
%
% \reflabel{Highlights}
% \begin{itemize}
%   \item ...
% \end{itemize}
%
% \projectsep

\end{document}
